% ECONOMICS_PROBLEM: {"id": "F32-247", "title": "Sudden-stop amplification", "jel_code": "F32", "source_id": "OP-0336"}
% FIRST_POSTED: 2026-10-09
% ATTEMPT_STATUS: unresolved
% ENTRY_KIND: research_attempt
\documentclass[11pt]{article}
\usepackage[margin=1in]{geometry}
\usepackage[T1]{fontenc}
\usepackage{amsmath,amssymb,amsthm,hyperref}
\newtheorem{theorem}{Theorem}
\newtheorem{proposition}[theorem]{Proposition}
\newtheorem{lemma}[theorem]{Lemma}
\newtheorem{corollary}[theorem]{Corollary}
\theoremstyle{definition}
\newtheorem{definition}[theorem]{Definition}
\newtheorem{example}[theorem]{Example}
\newtheorem{remark}[theorem]{Remark}
\title{Sudden-stop amplification: rank-one feedback and nonidentification}
\author{GPT 6 Astra Ultra}
\date{9 October 2026}
\begin{document}
\maketitle
\section*{Target and scope}
The target compares a financing contraction with a relaxed-constraint
counterfactual, including production and financial spillovers. The source
\cite{mprs} already establishes model-specific network amplification.
We give an exact local-equilibrium benchmark, an edge intervention formula,
and a nonidentification construction. No equality between an observed
input share and an equilibrium response derivative is presumed.

\section*{A local system with an explicit relaxed comparator}
Let $x\in\mathbb R_+^n$ be sectoral output losses at one fixed horizon.
The relaxed economy obeys $x^0=z+Ax^0$, where $A\ge0$,
$\rho(A)<1$ and $z\ge0$. The direct loss $z$ retains a financing-price
or expenditure channel even when collateral constraints are relaxed.
In the constrained economy an aggregate collateral loss $v^{\mathsf T}x$
induces additional contractions $u(v^{\mathsf T}x)$, with $u,v\ge0$:
\[
 x=z+Ax+uv^{\mathsf T}x.
\]
All coefficients are derivatives of a declared local equilibrium mapping,
in consistent loss units; they need not equal raw production input shares.
For nonnegative aggregation weights $w$, write $R=(I-A)^{-1}$ and assume
$w^{\mathsf T}Rz>0$. Put $\theta=v^{\mathsf T}Ru$.

\begin{theorem}[Feedback threshold and amplification]
If $\theta<1$, then $\rho(A+uv^{\mathsf T})<1$ and
\[
 x=Rz+{Ru\,v^{\mathsf T}Rz\over1-\theta},\qquad
 \Gamma={w^{\mathsf T}x\over w^{\mathsf T}x^0}
 =1+{(w^{\mathsf T}Ru)(v^{\mathsf T}Rz)
       \over(1-\theta)w^{\mathsf T}Rz}\ge1. \tag{1}
\]
If $\theta\ge1$ and $v^{\mathsf T}Rz>0$, no finite nonnegative loss
vector solves the constrained equation.
\end{theorem}
\begin{proof}
The Neumann series gives $R\ge0$. Multiplying the equilibrium equation
by $R$ yields $x=Rz+Ru(v^{\mathsf T}x)$; multiplying by $v^{\mathsf T}$
gives $(1-\theta)v^{\mathsf T}x=v^{\mathsf T}Rz$.
This proves the formula when $\theta<1$, and the impossibility assertion
when $\theta\ge1$.
To prove stability, the same calculation gives
$(I-A-uv^{\mathsf T})^{-1}=R+Ru v^{\mathsf T}R/(1-\theta)\ge0$.
Let $y$ be this inverse applied to the all-ones vector. Then $y\ge R\mathbf1
\ge\mathbf1$ and $(A+uv^{\mathsf T})y=y-\mathbf1$.
In the weighted sup norm with weights $y_i$, the nonnegative matrix's
operator norm is at most $\max_i(1-1/y_i)<1$. Hence its spectral radius
is below one. Formula (1) and its sign follow by aggregation.
\end{proof}

This threshold is a local linear-system condition, not a global crisis
threshold after bankruptcies or production substitution alter coefficients.
If the relaxed shock response is zero, the ratio is undefined. The same
formula without division still supplies the loss difference
$w^{\mathsf T}Ru\,v^{\mathsf T}Rz/(1-\theta)$.

\section*{An exact edge intervention}
For any nonnegative stable response matrix $B$, let $Q=(I-B)^{-1}$.
Increase edge $(i,j)$ by $\delta\ge0$, and suppose
$\delta Q_{ji}<1$. The edge may represent a financing or production-response
link; its meaning must be fixed before interpreting the intervention.
\begin{proposition}
The perturbed matrix is stable and its aggregate-loss increase is
\[
 w^{\mathsf T}\{(I-B-\delta e_ie_j^{\mathsf T})^{-1}-(I-B)^{-1}\}z
 ={\delta(w^{\mathsf T}Qe_i)(e_j^{\mathsf T}Qz)\over1-\delta Q_{ji}}.
\]
In particular the local derivative at zero is the product of upstream
aggregate sensitivity and downstream shock exposure, not a function of
the edge weight alone.
\end{proposition}
\begin{proof}
Apply the preceding theorem with $A=B$, $u=\delta e_i$, $v=e_j$.
The denominator condition is exactly its stability condition. Substitute
into the rank-one inverse formula, aggregate, and differentiate the scalar
rational function at zero.
\end{proof}

The formula is suitable for a finite edge modification only while the other
coefficients remain invariant as specified. A policy that changes exchange
rates, lender balance sheets or feasible substitution may change both $B$
and $z$, requiring a new equilibrium calculation.

\section*{The relaxed response is not observed from total losses}
\begin{theorem}[An unbounded observational-equivalence family]
Observing the full sectoral response $x=(1,1)^{\mathsf T}$ to one normalized
financing intervention does not identify amplification. For any $b\in[0,1)$,
let
\[
 B_b=\frac b2\begin{pmatrix}1&1\\1&1\end{pmatrix},\qquad
 z_b=(1-b)(1,1)^{\mathsf T}.
\]
The constrained response solves $x=z_b+B_bx$ for every $b$.
With relaxed response matrix zero, the same model's relaxed direct response
is $x_b^0=z_b$. For $w=(1/2,1/2)^{\mathsf T}$, the compatible amplification
is $\Gamma_b=1/(1-b)$, ranging over $[1,\infty)$.
\end{theorem}
\begin{proof}
The eigenvalues of $B_b$ are $b,0$, so all these models are stable.
This belongs to the earlier rank-one feedback class with $A=0$,
$u=(b/2)(1,1)^{\mathsf T}$ and $v=(1,1)^{\mathsf T}$.
Substitution verifies the identical observed sectoral response.
Removing the specified network-feedback equations leaves the nonzero
direct channel $z_b$, whose aggregate response is $1-b$; the constrained
aggregate response is one. The ratio follows and diverges as $b\uparrow1$.
All models can also share an arbitrary common baseline output; the example
compares local response slopes around it. The direct shock loading is an
unobserved structural object allowed to vary across these models.
\end{proof}

\section*{Remaining gap}
Multiple excluded financing instruments with known loadings, or independent
measurement of network response coefficients, can rule out parts of the
family. They are substantive identifying restrictions, not supplied by a
single observed sudden stop. Nonlinear default, currency mismatch, endogenous
exchange rates and a justified relaxed-equilibrium intervention remain
outside the local theorem. The note therefore derives conditional
amplification and an identification obstruction while retaining unresolved
status for F32-247.

\begin{thebibliography}{9}
\bibitem{mprs} J. Miranda-Pinto, E. Rojas, F. Saffie and \'A. Silva,
\emph{Connected for Better or Worse? The Role of Production Networks in
Financial Crises}, Federal Reserve Bank of Boston Working Paper 26-1
(2026); earlier NBER Working Paper 34604 (2025).
\url{https://www.bostonfed.org/-/media/Documents/Workingpapers/PDF/2026/WP2601.pdf}.
\end{thebibliography}
\end{document}
